\begin{trapbox}

	\begin{description}[style=nextline, leftmargin=2.8em, labelsep=0.5em, itemsep=0.35em]
		\item[\textbf{\textcolor{CTrap}{易错点一（有向比符号）}}]
			计算中直接用长度绝对值代入，忽略有向符号，导致结果错误.
		\item[\textbf{\textcolor{CTrap}{正确理解（正负判定）：}}]
			公式中的比必须是有向线段比，当点位于线段延长线上时，比值为负.
		\item[\textbf{\textcolor{CTrap}{错因分析（方向认知）：}}]
			学生习惯用正数比例，未建立有向概念.
	\end{description}

	\medskip
	\begin{description}[style=nextline, leftmargin=2.8em, labelsep=0.5em, itemsep=0.35em]
		\item[\textbf{\textcolor{CTrap}{易错点二（共点条件缺失）}}]
			只看到点分别在边上就直接套用定理，没有检查三线是否共点.
		\item[\textbf{\textcolor{CTrap}{正确理解（前提验证）：}}]
			塞瓦定理要求 $AD,BE,CF$ 三线共点，缺一不可.
		\item[\textbf{\textcolor{CTrap}{错因分析（忽视条件）：}}]
			认为点在边上自然满足定理，未严格验证.
	\end{description}

	\medskip
	\begin{description}[style=nextline, leftmargin=2.8em, labelsep=0.5em, itemsep=0.35em]
		\item[\textbf{\textcolor{CTrap}{易错点三（比例顺序出错）}}]
			代入时写成 $\frac{BD}{DC}\cdot\frac{EA}{CE}\cdot\frac{AF}{FB}=1$ 等错误顺序.
		\item[\textbf{\textcolor{CTrap}{正确理解（循环顺序）：}}]
			必须按照 $BD/DC,\; CE/EA,\; AF/FB$ 的顺序，每个比的分母对应下一个顶点.
		\item[\textbf{\textcolor{CTrap}{错因分析（记忆偏差）：}}]
			未理解结构是循环乘积，顺序影响结果.
	\end{description}

\end{trapbox}
